
Three experiments map directly to three research questions.

\textbf{RQ1 (Heterogeneity --- H-E1):} Does LLaMA-3.1-8B MLP layer-wise activation sparsity exhibit significant variation across 32 layers (CV > 0.3) and is the layer rank ordering stable across Alpaca vs. WikiText-103 (Kendall's tau\_calibration >= 0.6)?

\textbf{RQ2 (Cross-Distribution Stability --- H-M1):} Is the sparsity profile stable across four diverse calibration distributions? Gate: ICC(3,k) > 0.75 and all six pairwise Kendall's tau >= 0.6.

\textbf{RQ3 (Threshold Invariance --- H-M2):} Is the layer rank ordering invariant across epsilon in {0.001, 0.01, 0.05, 0.1}? Gate: CV > 0.3 for >= 3 of 4 epsilon values; maximum adjacent-pair Kendall's tau >= 0.7.

\subsubsection{4.1 Datasets}

\begin{table}[htbp]
\centering
\resizebox{\columnwidth}{!}{%
\begin{tabular}{llll}
\toprule
Dataset & HuggingFace Identifier & Domain & Samples Used \\
\midrule
Alpaca & tatsu-lab/alpaca & Instruction following & 512 (train split) \\
WikiText-103 & wikitext / wikitext-103-raw-v1 & General web text & 512 (test split, 512-token chunks) \\
SST-2 val & SetFit/sst2 & Sentiment classification & 512 (validation set) \\
MNLI val & nyu-mll/multi\_nli & Natural language inference & 512 (validation\_matched) \\
\bottomrule
\end{tabular}
}
\end{table}

For RQ1 (H-E1), Alpaca (lengths 128 and 512 tokens) and WikiText-103 are used. For RQ2 (H-M1), all four distributions are used. For RQ3 (H-M2), Alpaca and WikiText-103 are used across all four epsilon values.

\subsubsection{4.2 Evaluation Metrics}

\textbf{RQ1:} CV of sparsity across 32 layers (Alpaca, length 512, epsilon = 0.01); Kendall's tau\_calibration between Alpaca and WikiText-103 sparsity rankings; Kendall's tau\_length between 128-token and 512-token Alpaca profiles.

\textbf{RQ2:} ICC(3,k) computed across all four distributions; all 6 pairwise Kendall's tau values.

\textbf{RQ3:} CV per epsilon (pass rate >= 3/4 required); all 6 cross-epsilon pairwise Kendall's tau values; maximum adjacent-pair tau.

All Kendall's tau values are computed with scipy.stats.kendalltau (variant='b', two-tailed p-values). ICC is computed with pingouin 0.6.1.

